\section*{Publisher term/frame and measure conventions}
Let $\Sigma$ be a set of function signatures, $V$ be a set of variables, and $N$ be a set of names. Each function signature $f\in\Sigma$ has arity of some non-negative integer $n\ge0$.
The set $T(\Sigma,X)$ of $\Sigma$-terms is defined as the smallest set such that $X\subseteq T(\Sigma,X)$, and if $f\in\Sigma^{(n)}$ and $t_1,\ldots,t_n\in T(\Sigma,X)$, then $f(t_1,\ldots,t_n)\in T(\Sigma,X)$. Unless explicitly specified, $X=V\cup N$.
A term is ground if it has no variable. A substitution $\sigma:V\to T(\Sigma,X)$ differs from the identity on only finitely many variables; $t\sigma$ acts recursively on variables and function arguments, leaving names fixed. We use $\operatorname{dom}(\sigma)$ and $\operatorname{ran}(\sigma)$ as in publisher section2.1, and $\sigma\circ\tau(x)=\sigma(\tau(x))$.
Given a term $t$, $\operatorname{pos}(\alpha,t)=\{p\in\operatorname{pos}(t):t(p)=\alpha\}$ and $n_\alpha(t)=|\operatorname{pos}(\alpha,t)|$. We write $t|_p$ for the subterm at position $p$, and $\epsilon$ for the root position.
The equational theory $\approx_E$ is the smallest congruence on $T(\Sigma,X)$ containing all equations $l\approx r\in E$ that is closed under substitutions and contexts. The associated rewrite rules are $R_E=\{l\to r:(l\approx r)\in E\}$; $s\to_E t$ iff $s|_p=l\sigma$ and $t=s[r\sigma]_p$ for some rule, position and substitution. $\leftrightarrow_E^*$ denotes the symmetric-transitive closure, and $\leftrightarrow_E^\epsilon$ denotes rewriting at the root.
A frame is $\varphi=\nu\tilde n.\sigma=\nu\tilde n.\{x_1\mapsto t_1,\ldots,x_n\mapsto t_n\}$ for a ground substitution and a set of names $\tilde n$.
A term $t\in T(\Sigma,X)$ is a recipe for $\varphi$ if $\operatorname{name}(t)\cap\tilde n=\varnothing$.
\paragraph{Definition2 (Deduction).} A term $t$ is deducible from $\varphi$ if there exists a recipe $r$ for $\varphi$ such that $r\sigma\to_E^*t$. We denote this by $\varphi\vdash_E t$.
\paragraph{Definition4 (Static Equivalence).} For $\varphi_1=\nu\tilde n.\sigma_1$ and $\varphi_2=\nu\tilde n.\sigma_2$, $\varphi_1\equiv_E\varphi_2$ if for every terms $r,s$ such that $(\operatorname{name}(r)\cup\operatorname{name}(s))\cap\tilde n=\varnothing$, $r\sigma_1\approx_Es\sigma_1$ iff $r\sigma_2\approx_Es\sigma_2$.
\paragraph{Definition11 ($w$-Measure).} Given $w:\Sigma\cup N\to\mathbb N_{>0}$, the $w$-measure $\phi:T(\Sigma,V\cup N)\to\mathbb N_{>0}$ satisfies
\[
\phi(f(t_1,\ldots,t_n))=w(f)+\sum_{i=1}^n\phi(t_i),\quad \phi(x)=1\ (x\in V),\quad \phi(n)=w(n)\ (n\in N).
\]
\paragraph{Definition12 (Measure-Invariant Equational Theory).} The equation $l\approx r$ is $w$-invariant if $n_x(l)=n_x(r)$ for every $x\in V$, and
\[
\sum_{f\in\Sigma}w(f)n_f(l)+\sum_{a\in N}w(a)n_a(l)=\sum_{f\in\Sigma}w(f)n_f(r)+\sum_{a\in N}w(a)n_a(r).
\]
An equational theory $E$ is measure-invariant if there exists $w$ such that every equation in $E$ is $w$-invariant. A TRS $R$ is measure-invariant if $R=R_E$ for such an $E$.
\paragraph{Definition17 (Separate Measure-Invariant (SMI) Equational Theory).} $E$ is SMI if $E$ is measure-invariant and for any $l\approx r\in E$, $l$ and $r$ are rooted by the same function symbol that does not occur elsewhere in $E$.

\section*{Complete publisher measure-invariant deduction proof and AppendixA continuation}
\paragraph{Proposition42.} Let $\sigma=\{x_1\mapsto t_1,\ldots,x_n\mapsto t_n\}$. For any $s$ with $\operatorname{var}(s)\subseteq\{x_1,\ldots,x_n\}$,
\[
\phi(s\sigma)=\sum_{f\in\Sigma}w(f)n_f(s)+\sum_{a\in N}w(a)n_a(s)+\sum_{i=1}^nn_{x_i}(s)\phi(t_i).
\]
\paragraph{Lemma43.} For any $s,\sigma$, $\phi(s\sigma)\ge\phi(s)$.
\begin{proof}This is by Proposition42 and that $\phi(t)\ge1=\phi(x)$ for every term $t$.\end{proof}
\paragraph{Lemma44.} If $R$ is measure-invariant by $w$, then for any $l\approx r\in R$ and any $w$-measure $\phi$, $\phi(l\sigma)=\phi(r\sigma)$ for all substitution $\sigma$.
\begin{proof}
Suppose $l\approx r$ is measure-invariant. Let $\sigma=\{x_1\mapsto t_1,\ldots,x_n\mapsto t_n\}$. Without loss of generality $\operatorname{var}(l)=\operatorname{var}(r)\subseteq\{x_1,\ldots,x_n\}$ and by measure-invariant property, $x_i$ occurring in $l$ and $r$, note that $n_{x_i}(l)=n_{x_i}(r)$.
Since $x_i$ are substituted to $t_i$ in both $l$ and $r$, then $n_{x_i}(l)\phi(t_i)=n_{x_i}(r)\phi(t_i)$.
Moreover, by the second condition of measure-invariance of $l\approx r$, we have
\[
\sum_{f\in\Sigma}w(f)n_f(l)+\sum_{a\in N}w(a)n_a(l)=\sum_{f\in\Sigma}w(f)n_f(l)+\sum_{a\in N}w(a)n_a(l),
\]
and hence $\phi(l\sigma)=\phi(r\sigma)$.
\end{proof}
\paragraph{Lemma14.} Let $R$ be measure-invariant by $w$. If $s\to_Rt$, then $\phi(s)=\phi(t)$ where $\phi$ is a $w$-measure.
\begin{proof}
Let $p$ be a position and $\sigma$ be a substitution such that $s|_p=l\sigma$ and $t=s[r\sigma]_p$ for some $l\approx r\in R$.
Then $\phi(t)=\phi(s)-\phi(l\sigma)+\phi(r\sigma)=\phi(s)$.
\end{proof}
\paragraph{Theorem15.} Deduction is decidable for any measure-invariant equational theory.
\begin{proof}
Let $\varphi=\nu\tilde n.\sigma$ be a frame and $t$ be a (target) term. Then, $t$ is deducible iff there exists a term $s$ with $n(s)\cap\tilde n=\varnothing$ such that $s\varphi\to^*t$.
Note that $s\varphi\to s_1\to s_2\to\cdots\to s_n=t$.
Then, by Lemma14, $\phi(s\sigma)=\phi(s_1)=\cdots=\phi(s_n)=\phi(t)$.
Moreover, for every substitution $\sigma$, we have $\phi(s\sigma)\ge\phi(s)$ and so $\phi(s)\le\phi(t)$.
Hence, there are only finitely many terms $s$ to check (up to renaming: if the term $s$ uses names from $N\setminus(n(E)\cup\tilde n)$, the infinite choices of the names will be indifferent, and can be fixed to at most $\phi(s)$ of them).
\end{proof}

\paragraph{Theorem19.} Deduction and static equivalence problem for SMI theory is decidable.

\section*{Complete publisher AppendixB: Decidability of Static Equivalence for SMI Theory}
The proof follows a standard saturation-based approach to decidability, adapting the technique of [1] for subterm-convergent theories, which has since been widely applied to establish decidability of static equivalence in related settings [9,18,21].
Given a frame $\varphi=\nu\tilde n.\sigma$, we construct a finite saturation set of deducible terms, each equipped with a corresponding recipe. From this set, we extract a finite characteristic set of equalities whose terms are bounded in size by the measure of the theory. We then show that two frames are statically equivalent if and only if each satisfies the characteristic equalities of the other.
For an SMI equational theory $E$ with weight $w$ and measure $\phi$, define $\phi(E)=\max_{l\approx r\in E}\phi(l)$. For any term $t$, let $\operatorname{St}(t)$ denote its subterms; $\operatorname{St}(\sigma)=\bigcup_{t\in\operatorname{ran}(\sigma)}\operatorname{St}(t)$.
\paragraph{Definition45 (Extended Set of Subterms).} Let
\[
\operatorname{ESt}(\sigma)=\operatorname{St}(\sigma)\cup\bigcup_{t\in ET}\operatorname{St}(t),
\]
where $ET$ is the set of terms $t$ such that $s[d_1,\ldots,d_n]\leftrightarrow_E^\epsilon t$ for some context $s$ such that $\phi(s)\le\phi(E)$, $fn(s)\cap\tilde n=\varnothing$, and $d_1,\ldots,d_n\in\operatorname{St}(\sigma)$.
Recall that for SMI theory $E$, $s\leftrightarrow_Et$ iff $\phi(s)=\phi(t)$, so the amount of terms that $s$ can reach using $\leftrightarrow_E$ is finite. Hence, thet set $ET$ is finite. It follows that the set $\operatorname{ESt}(\sigma)$ is finite as well, given that $\sigma$ is finite.
\paragraph{Definition46 (Saturation Set).} For any frame $\varphi=\nu\tilde n.\sigma$, the set $\operatorname{sat}_E(\varphi)$ is the smallest set $D$ such that:
\begin{enumerate}
\item $\operatorname{ran}(\sigma)\subseteq D$;
\item if $t_1,\ldots,t_n\in D$ and $f(t_1,\ldots,t_n)\in\operatorname{ESt}(\sigma)$, then $f(t_1,\ldots,t_n)\in D$;
\item if $t'\in D$, $t\in\operatorname{ESt}(\sigma)$, $t'\approx_Et$, then $t\in D$.
\end{enumerate}
Note that $\operatorname{sat}_E(\varphi)$ is finite. Rule2 and3 guarantee this since it only accepts $f(t_1,\ldots,t_n)$ that belongs to $\operatorname{ESt}(\sigma)$, which is finite.
Note also that if $t\in\operatorname{sat}_E(\varphi)$, then $t$ is deducible from $\varphi$, that is there exists a recipe $\zeta_t$ such that $\zeta_t\varphi\approx_Et$. We can easily prove this by induction on the construction of $\operatorname{sat}_E(\varphi)$.
If the term $t$ is included in the set by using rule1, the recipe is just the domain variable in $\varphi$ corresponding to $t$.
If $t$ is included by rule2, then $\zeta_t=f(\zeta_1,\ldots,\zeta_n)$ where $\zeta_i$ is the recipe for $t_i$, for $i=1,\ldots,n$.
If $t$ is included by rule3, it uses the same recipe as $t'$.
\paragraph{Definition47 (Completion of a Frame).} $\sigma^*=\sigma\cup\{\chi_u\mapsto u:u\in\operatorname{sat}_E(\varphi)\setminus\operatorname{ran}(\sigma)\}$, where $\chi_u$ is a fresh variable for any such $u$, and $\varphi^*$ denotes $\nu\tilde n.\sigma^*$.
Given a recipe $\zeta_u$ for each such $u$, the substitution $\{\chi_u\mapsto\zeta_u:u\in\operatorname{sat}_E(\varphi)\setminus\operatorname{ran}(\sigma)\}$ is called a recipe substitution of $\varphi$ and is denoted $\zeta_\varphi$.
\paragraph{Definition48 (Characteristic Set of Equalities).} Let $B_t=\{t:\phi(t)\le\phi(E)\}$. The set $\operatorname{Eq}_B(\varphi)$ is the set of equalities $t\zeta_\varphi=t'\zeta_\varphi$ such that $(t\zeta_\varphi\approx_Et'\zeta_\varphi)\varphi$ and $t,t'\in B_t$.
\paragraph{Lemma49.} Let $E$ be an SMI theory. For any $s,t$ satisfying the name restriction, if $s\varphi^*\approx_Et\varphi^*$ and $\psi\models\operatorname{Eq}_B(\varphi)$, then $(s\zeta_\varphi)\psi\approx_E(t\zeta_\varphi)\psi$.
\begin{proof}
We prove by induction on the size of $\max(\phi(s),\phi(t),\phi(E))$.
Base case: if $\max(\phi(s),\phi(t),\phi(E))=\phi(E)$, then it is true by definition of $\operatorname{Eq}_B$.
Inductive step: Let $k>\phi(E)$ and suppose that the statement is true for any $s,t$ such that $\max(\phi(s),\phi(t),\phi(E))\le k$.
Consider $s,t$ are terms such that $\max(\phi(s),\phi(t),\phi(E))=k$. Let
\[
s=f(s_1,\ldots,s_n),\qquad t=f(t_1,\ldots,t_n)
\]
and $s\varphi^*\approx_Et\varphi^*$.
This means that $f(s_1\varphi^*,\ldots,s_n\varphi^*)\leftrightarrow_E^*f(t_1\varphi^*,\ldots,s_n\varphi^*)$.
If $s_i\varphi^*\approx_Et_i\varphi^*$ for every $i=1,\ldots,n$, then by IH, we have $(s_i\zeta_\varphi)\psi\approx_E(t_i\zeta_\varphi)\psi$ and hence, $f(s_1,\ldots,s_n)\zeta_\varphi\psi\approx_E f(t_1,\ldots,t_n)\zeta_\varphi\psi$.
Suppose not, then there exist $f(l_1,\ldots,l_n)\approx f(r_1,\ldots,r_n)\in E$ and a substitution $\rho$ such that
\[
s\varphi^*=f(s_1\varphi^*,\ldots,s_n\varphi^*)=f(l_1,\ldots,l_n)\rho,
\qquad t\varphi^*=f(t_1\varphi^*,\ldots,t_n\varphi^*)=f(r_1,\ldots,r_n)\rho.
\]
This means that $s_i\varphi^*\approx_El_i\rho$ and $t_i\varphi^*\approx_Er_i\rho$, for $i=1,\ldots,n$, which is basically a syntactic unification problem instance, since all function symbols in $l_i$ and $r_i$ are only inner constructors in $E$.
Consider all variables $x$ occuring in $l_1,\ldots,l_n$ (same as in $r_1,\ldots,r_n$). Because of the equations $s_i\varphi^*=l_i\rho$ and $t_i\varphi^*=r_i\rho$, there are two possible cases of what $x\rho$ is:
\begin{enumerate}
\item $x\rho=u[w_1,\ldots,w_m]$ for some $w_1,\ldots,w_m\in\operatorname{ran}(\varphi^*)$.
\item $x\rho=u|_p$ for some $u\in\operatorname{ran}(\varphi^*)$ and $p\ne\epsilon$ such that any subterm $u|_q$ for $\epsilon\le q<p$ is rooted by an inner constructor. In this case, note that $u|_p\in\operatorname{ran}(\varphi^*)$ and so there exists a variable $y$ such that $y\varphi^*=u|_p$.
\end{enumerate}
We define another substitution $\mu$ such that if it is the first case, we let $x\mu=u[v_1,\ldots,v_m]$ where $v_j\varphi^*=w_j$ for $j=1,\ldots,m$ and if it is the second case, we define $x\mu=y$.
Using this definition of $\mu$, we have $\rho=\varphi^*\circ\mu$.
To show this, note that in case1, we have $(x\mu)\varphi^*=u[v_1,\ldots,v_j]\varphi^*=u[v_1\varphi^*,\ldots,v_m\varphi^*]=u[w_1,\ldots,w_m]=\rho(x)$.
In the second case, $\varphi^*(\mu(x))=\varphi^*(y)=u|_p=\rho(x)$, as desired.
Now, note that for every $i=1,\ldots,n$, we have that $k\ge\phi(s)>\phi(s_i\varphi^*)=\phi(l_i\rho)=\phi(\varphi^*(\mu(l_i)))\ge\phi(\mu(l_i))$, and similarly, $k>\phi(\mu(r_i))$.
By IH, we have also for every $i=1,\ldots,n$, since $s_i\varphi^*\approx_E(l_i\mu)\varphi^*$ and $t_i\varphi^*\approx_E(r_i\mu)\varphi^*$, implying that
\[
(s_i\zeta_\varphi)\psi\approx_E(l_i\mu\zeta_\varphi)\psi,\qquad (t_i\zeta_\varphi)\psi\approx_E(r_i\mu\zeta_\varphi)\psi\quad(i=1,\ldots,n).
\]
This, we have
\begin{align*}
(s\zeta_\varphi)\psi&=f((s_1\zeta_\varphi)\psi,\ldots,(s_n\zeta_\varphi)\psi)\\
&\approx_Ef((l_1\mu\zeta_\varphi)\psi,\ldots,(l_n\mu\zeta_\varphi)\psi)\\
&\approx_Ef((r_1\mu\zeta_\varphi)\psi,\ldots,(r_n\mu\zeta_\varphi)\psi)\\
&\approx_Ef((t_1\zeta_\varphi)\psi,\ldots,(t_n\zeta_\varphi)\psi)=(t\zeta_\varphi)\psi.
\end{align*}
so the inductive case is completed.
\end{proof}
\paragraph{Lemma50.} For any frames $\varphi,\psi$, $\varphi\approx_E\psi$ iff $\psi\models\operatorname{Eq}_B(\varphi)$ and $\varphi\models\operatorname{Eq}_B(\psi)$.
\begin{proof}
The only if direction is clear. For the if direction, let $\operatorname{Eq}(\varphi)$ be the set of all equalities $s\zeta_\varphi=t\zeta_\varphi$ such that $(s\zeta_\varphi\approx_Et\zeta_\varphi)\varphi$.
Take any $(s\zeta_\varphi\approx_Et\zeta_\varphi)\in\operatorname{Eq}(\varphi)$.
By definition of $\zeta_\varphi$ and $\varphi^*$, we have $(s\zeta_\varphi)\varphi\approx_Es\varphi^*$ and $(t\zeta_\varphi)\varphi\approx_Et\varphi^*$.
Since $(s\zeta_\varphi)\varphi\approx_E(t\zeta_\varphi)\varphi$, we obtain $s\varphi^*\approx_Et\varphi^*$.
Then, by Lemma49, we get $(s\zeta_\varphi)\psi\approx_E(t\zeta_\varphi)\psi$, which means that $\psi\models\operatorname{Eq}(\varphi)$.
Similarly, we show that $\varphi\models\operatorname{Eq}(\psi)$. Then, the conclusion follows.
\end{proof}
\paragraph{Proof of Theorem19.}
The deduction is decidable because $E$ is measure-invariant (by Lemma15). Static equivalence is decidable by and Lemma49 and Lemma50.

